\lecture{2}{Thu 11 Jan 2024 13:32}{Homology}

\subsection{Homology}

Elements of $Z_n$ are called $n$-cycles, $B_n$ the $n$-boundaries, and $H_n$ the homology classes.

If two elements $c, c'$ in $Z_n$ are inside the same equivalence class, $[c] = [c'] \in H_n$, then we say $c$ and $c'$ are homologous.

The chain complex looks like:

\[
	\ldots \to S_1(X) \to S_0(X) \to 0 \to 0 \to  \ldots
.\] 

\begin{notation}
	Let $c_x^n$ be the map
	
	\begin{align*}
		c_x^n: \Delta^n &\longrightarrow X \\
		\nu &\longmapsto c_x^n(\nu) = x \in X
	.\end{align*}
\end{notation}

Thus
\[
	S_0(X) = \mathbb{Z} \{c_x^0 | x \in X\} 
.\] 

\[
	S_1(X) = \mathbb{Z} \{ \text{paths in }X\} 
.\] 

\begin{example}(Understanding the Homologies)
	
	\begin{itemize}
		\item The zeroth homology.

$B_0(X)$ consists of boundaries of points connected by a path in $S_1(X)$.

$H_0(X)$ is the free abelian group generated by connected components of $X$, since for any two points in $S_0(X)$, they are the same in the quotient $B_0(X)$ iff they are path-connected.


	\item The first homology. 
		Take $\sigma \in \text{Sin}_1(X)$, then $d \sigma = \sigma|_{[1]} - \sigma|_{[0]}$. Thus $\sigma \in Z_1(X)$ iff $\sigma$ represents a loop in $X$. For $\sigma \in S_1(X)$, $d \sigma = 0$ iff they together form a directed closed loop. 

		Define $\overline{\sigma} $ to be the reversion of $\sigma$ : $\overline{\sigma} (t) = \sigma(1 - t), t \in [0,1]$. 

		Claim that $\sigma + \overline{\sigma} \in B_1(X)$. To verify, define an affine projection $\pi: \Delta_2 \to \Delta_1$. Let $\pi(0) = 0, \pi(2) = 0, \pi(1) = 1$. Now $\sigma \circ \pi \in S_2(X)$. Compute
		\[
			d(\sigma \circ \pi)
			= \sigma \circ \pi|_{[1 2]} - \sigma \circ \pi |_{[0 2]} + \sigma \circ \pi _{[0 1]}
			= \overline{\sigma}  - c^1_{\sigma(0)} + \sigma
		.\] 
		If consider $c^2_{\sigma(0)}$, claim that  $d c^2_{\sigma(0)} = c^1_{\sigma(0)}$.
		Thus,
		\[
			\sigma + \overline{\sigma} = d\left( \sigma \circ \pi + c^2_{\sigma(0)} \right) 
		.\] 
		Therefore (check), if $\sigma \in Z_1(X)$ is a loop, then $[\overline{\sigma} ] = - [\sigma]$.

		For HW, for path connected $X$, check $H_1(X)$ is the Abelinization of $\pi_1(X)$.
	\end{itemize}
\end{example}

\begin{example}(Computing Homologies)
	
	\begin{itemize}
		\item Topological space with one point: $S_n(\cdot )$. We have $S_n(\cdot ) = \mathbb{Z} [c_\cdot ^n]$. We see
\begin{align*}
	d(c_*^n) &= \sum_{i = 0}^n (-1)^i c_*^n |_{[0 \ldots \hat{i}  \ldots n]}\\
	= \begin{cases}
		c_*^{n - 1}  & n \text{ even}\\
		0 & n \text{ odd}
	\end{cases}
.\end{align*}
We have
\[
	H_n(\cdot ) = \begin{cases}
		0 & n \neq 0 \\
		\mathbb{Z} & n = 0
	\end{cases}
.\] 
	\end{itemize}
\end{example}

For $f : X \to  Y$ continuous, we can generate a map
\begin{align*}
	f_*: \text{Sin}_n(X) &\longrightarrow \text{Sin}_n(Y) \\
	\sigma &\longmapsto f_*(\sigma) = f \circ \sigma
.\end{align*}
By extending linearly, we have
\begin{align*}
	f_*: S_n(X) &\longrightarrow S_n(Y) \\
	\sum n_\sigma \sigma &\longmapsto \sum n_\sigma f \circ \sigma
.\end{align*}

% https://quiver.local:8080/#q=WzAsNCxbMCwwLCJTX24oWCkiXSxbMiwwLCJTX3tuLTF9KFgpIl0sWzAsMiwiU19uKFkpIl0sWzIsMiwiU197bi0xfShZKSJdLFswLDEsImRfWCJdLFswLDIsImZfKiIsMl0sWzIsMywiZF9ZIiwyXSxbMSwzLCJmXyoiXV0=
\[\begin{tikzcd}
	{S_n(X)} && {S_{n-1}(X)} \\
	\\
	{S_n(Y)} && {S_{n-1}(Y)}
	\arrow["{d_X}", from=1-1, to=1-3]
	\arrow["{f_*}"', from=1-1, to=3-1]
	\arrow["{d_Y}"', from=3-1, to=3-3]
	\arrow["{f_*}", from=1-3, to=3-3]
\end{tikzcd}\]

We claim this diagram commutes. It suffices to check this on a basis element of $S_n(X)$. Take $\sigma: \Delta_n \to  X$. 
\begin{align*}
	f_* \circ d_X(\sigma) 
	&= f_*\left( \sum_{i=0}^{n} (-1)^i \sigma \circ d^i \right) \\
	&=  \sum_{i=0}^{n} (-1)^i f \circ \sigma \circ d^i  \\
	d_Y \circ f_* 
	&= d_Y(f \circ \sigma) = \sum_{i=0}^{n} (-1)^i f\circ \sigma \circ d^i 
.\end{align*}

\begin{definition}[Chain Map]
	\label{defn:ChainMap}
	Let $C_*, D_*$ be chain complexes, a \textit{chain map} $\phi: C_* \to D_*$ is $\{\phi_n: C_n \to  D_n\} _{n \in \mathbb{Z}}$ such that

% https://quiver.local:8080/#q=WzAsOCxbMiwwLCJDX24iXSxbNCwwLCJDX3tuLTF9Il0sWzIsMiwiRF9uIl0sWzQsMiwiRF97bi0xfSJdLFswLDIsIlxcY2RvdHMiXSxbNiwyLCJcXGNkb3RzIl0sWzYsMCwiXFxjZG90cyJdLFswLDAsIlxcY2RvdHMiXSxbMCwyLCJcXHBoaV9uIl0sWzEsMywiXFxwaGlfe24tMX0iXSxbMCwxLCJkX0MiLDFdLFsyLDMsImRfRCIsMV0sWzQsMl0sWzMsNV0sWzEsNl0sWzcsMF1d
\[\begin{tikzcd}
	\cdots && {C_n} && {C_{n-1}} && \cdots \\
	\\
	\cdots && {D_n} && {D_{n-1}} && \cdots
	\arrow["{\phi_n}", from=1-3, to=3-3]
	\arrow["{\phi_{n-1}}", from=1-5, to=3-5]
	\arrow["{d_C}"{description}, from=1-3, to=1-5]
	\arrow["{d_D}"{description}, from=3-3, to=3-5]
	\arrow[from=3-1, to=3-3]
	\arrow[from=3-5, to=3-7]
	\arrow[from=1-5, to=1-7]
	\arrow[from=1-1, to=1-3]
\end{tikzcd}\]	
commutes.
\end{definition}

\begin{claim}
	$\phi_n$ induces a map $H_n (C_*) \to H_n(D_*)$
\end{claim}
\begin{proof}
	We check that $\phi_n$ maps $B_n(C_*)$ to $B_n(D_*)$, and maps $Z_n(C_*) $ to $Z_n(D_*)$.

	Chase the diagram. If some $\sigma \in B_{n - 1}(C_*)$, then there exists $\tau \in C_n$ such that $\sigma = d_c \tau$. Then $\phi_{n - 1} \sigma = \phi_{n - 1} d_c \tau = d_D (\phi_n \tau) \in B_{n - 1}(D_*)$.
\end{proof}

Thus we have the map
\begin{align*}
	f_*: H_n(X) &\longrightarrow H_n(Y) \\
	\left[ \sum n_\sigma \sigma \right]  &\longmapsto \left[ \sum n_\sigma f \circ \sigma \right] 
.\end{align*}









